\newpage

% =========================================================
% ВАРИАНТ 5
% =========================================================

\section*{Вариант 5}

Биссектриса тупого угла параллелограмма делит противоположную сторону
в отношении \(3:4\), считая от вершины острого угла.
Найдите большую сторону параллелограмма, если его периметр равен \(33\).

\begin{center}
\begin{tikzpicture}[
  scale=0.72,
  line cap=round,
  line join=round,
  every node/.style={font=\small}
]
  \coordinate (A) at (0,0);
  \coordinate (B) at (7,0);
  \coordinate (E) at (3,0);
  \coordinate (D) at (0.8,2.891);
  \coordinate (C) at (7.8,2.891);

  % Параллелограмм
  \draw[line width=0.9pt]
    (A)--(B)--(C)--(D)--cycle;

  % Биссектриса тупого угла D
  \draw[line width=0.9pt] (D)--(E);

  % Обозначение равных углов
  \pic[
    draw,
    line width=0.65pt,
    angle radius=5mm
  ] {angle=A--D--E};

  \pic[
    draw,
    line width=0.65pt,
    angle radius=5mm
  ] {angle=E--D--C};

  % Подписи точек
  \node[below left]  at (A) {$A$};
  \node[below right] at (B) {$B$};
  \node[above right] at (C) {$C$};
  \node[above left]  at (D) {$D$};
  \node[below]       at (E) {$E$};

  % Отношение частей стороны
  \node[below=4mm] at ($(A)!0.5!(E)$) {$3$};
  \node[below=4mm] at ($(E)!0.5!(B)$) {$4$};
\end{tikzpicture}
\end{center}

\[
\boxed{\text{Ответ: }11{,}55}
\]


% =========================================================
% ВАРИАНТ 6
% =========================================================

\section*{Вариант 6}

Точка пересечения биссектрис двух углов параллелограмма,
прилежащих к одной стороне, принадлежит противоположной стороне.
Меньшая сторона параллелограмма равна \(6\).
Найдите его большую сторону.

\begin{center}
\begin{tikzpicture}[
  scale=0.78,
  line cap=round,
  line join=round,
  every node/.style={font=\small}
]
  \coordinate (A) at (0,0);
  \coordinate (B) at (6,0);
  \coordinate (D) at (1,2.828);
  \coordinate (C) at (7,2.828);

  % Точка пересечения биссектрис
  \coordinate (E) at (3,0);

  % Параллелограмм
  \draw[line width=0.9pt]
    (A)--(B)--(C)--(D)--cycle;

  % Биссектрисы углов D и C
  \draw[line width=0.9pt] (D)--(E)--(C);

  % Биссектриса угла D
  \pic[
    draw,
    line width=0.65pt,
    angle radius=4.8mm
  ] {angle=A--D--E};

  \pic[
    draw,
    line width=0.65pt,
    angle radius=4.8mm
  ] {angle=E--D--C};

  % Биссектриса угла C
  \pic[
    draw,
    line width=0.65pt,
    angle radius=4.8mm
  ] {angle=D--C--E};

  \pic[
    draw,
    line width=0.65pt,
    angle radius=4.8mm
  ] {angle=E--C--B};

  % Подписи
  \node[below left]  at (A) {$A$};
  \node[below right] at (B) {$B$};
  \node[above right] at (C) {$C$};
  \node[above left]  at (D) {$D$};
  \node[below]       at (E) {$E$};
\end{tikzpicture}
\end{center}

\[
\boxed{\text{Ответ: }12}
\]


% =========================================================
% ВАРИАНТ 7
% =========================================================

\newpage

\section*{Вариант 7}

Острый угол прямоугольного треугольника равен \(50^\circ\).
Найдите угол между высотой \(CH\) и медианой \(CM\),
проведёнными из вершины прямого угла.
Ответ дайте в градусах.

\begin{center}
\begin{tikzpicture}[
  scale=0.9,
  line cap=round,
  line join=round,
  every node/.style={font=\small}
]
  \coordinate (A) at (0,0);
  \coordinate (B) at (6,0);

  % M — середина гипотенузы AB
  \coordinate (M) at (3,0);

  % H — основание высоты CH
  \coordinate (H) at (3.521,0);

  % Вершина прямого угла
  \coordinate (C) at (3.521,2.954);

  % Гипотенуза с отметками равных частей
  \draw[line width=0.9pt,equal segment] (A)--(M);
  \draw[line width=0.9pt,equal segment] (M)--(B);

  % Катеты
  \draw[line width=0.9pt] (A)--(C)--(B);

  % Медиана и высота
  \draw[line width=0.85pt] (C)--(M);
  \draw[line width=0.85pt] (C)--(H);

  % Прямой угол при C
  \pic[
    draw,
    line width=0.65pt,
    angle radius=3.5mm
  ] {right angle=A--C--B};

  % CH перпендикулярна AB
  \pic[
    draw,
    line width=0.65pt,
    angle radius=3.2mm
  ] {right angle=C--H--B};

  % Угол B = 50 градусов
  \pic[
    draw,
    line width=0.65pt,
    angle radius=7mm
  ] {angle=C--B--A};

  \node at (5.15,0.48) {$50^\circ$};

  % Подписи
  \node[below left]  at (A) {$A$};
  \node[below right] at (B) {$B$};
  \node[above]       at (C) {$C$};
  \node[below]       at (M) {$M$};
  \node[below]       at (H) {$H$};
\end{tikzpicture}
\end{center}

\[
\boxed{\text{Ответ: }10}
\]


% =========================================================
% ВАРИАНТ 8
% =========================================================

\section*{Вариант 8}

Угол между биссектрисой \(CD\) и медианой \(CM\),
проведёнными из вершины прямого угла \(C\) треугольника \(ABC\),
равен \(10^\circ\).
Найдите меньший угол этого треугольника.
Ответ дайте в градусах.

\begin{center}
\begin{tikzpicture}[
  scale=0.9,
  line cap=round,
  line join=round,
  every node/.style={font=\small}
]
  \coordinate (A) at (0,0);
  \coordinate (B) at (6,0);

  % M — середина гипотенузы
  \coordinate (M) at (3,0);

  % D — точка пересечения биссектрисы с гипотенузой
  \coordinate (D) at (3.528,0);

  % Вершина прямого угла
  \coordinate (C) at (4.026,2.819);

  % Гипотенуза с отметками равных частей
  \draw[line width=0.9pt,equal segment] (A)--(M);
  \draw[line width=0.9pt,equal segment] (M)--(B);

  % Катеты
  \draw[line width=0.9pt] (A)--(C)--(B);

  % Медиана и биссектриса
  \draw[line width=0.85pt] (C)--(M);
  \draw[line width=0.85pt] (C)--(D);

  % Прямой угол
  \pic[
    draw,
    line width=0.65pt,
    angle radius=3.5mm
  ] {right angle=A--C--B};

  % Обозначения равных частей угла C
  \pic[
    draw,
    line width=0.65pt,
    angle radius=7mm
  ] {angle=A--C--D};

  \pic[
    draw,
    line width=0.65pt,
    angle radius=7mm
  ] {angle=D--C--B};

  % Угол между медианой и биссектрисой
  \pic[
    draw,
    line width=0.65pt,
    angle radius=4.5mm
  ] {angle=M--C--D};

  \node at (3.68,2.10) {$10^\circ$};

  % Подписи
  \node[below left]  at (A) {$A$};
  \node[below right] at (B) {$B$};
  \node[above]       at (C) {$C$};
  \node[below]       at (M) {$M$};
  \node[below]       at (D) {$D$};
\end{tikzpicture}
\end{center}

\[
\boxed{\text{Ответ: }35}
\]

\end{document}